%--------------------------------------------------------------------
%
%   Version: 3.25        (Koma-Skript, documentclass{scrbook})
%   Datum: 01.06.2018
%   Autor: Johannes Bier
%
%--------------------------------------------------------------------
% die Auflistung für das Abkürzungsverzeichnis 
%--------------------------------------------------------------------

\newacronym{UVLO}{UVLO}{Undervoltage-lockout}
\newacronym{GIT}{GIT}{Gate Injection Transistoren}
\newacronym{NPT}{NPT}{Non-Punch-Through}
\newacronym{HEMT}{HEMT}{high-electron-mobility-transistor}
\newacronym{E-HEMT}{E-HEMT}{enhancement-mode high-electron-mobility-transistor}

\newacronym{WBG}{WBG}{wide-bandgap}
\newglossaryentry{BNC}{
	name={BNC},
	description={Bayonet-Neill-Concelman Bajonettverschluss},
	first={\glsentrydesc{BNC} (\glsentrytext{BNC})},
	plural={BNC},
	symbol={BNC},
	descriptionplural={Bayonet-Neill-Concelman Bajonettverschlüsse},
	firstplural={\glsentrydescplural{BNC} (\glsentryplural{BNC})}
}
\newacronym{LWL}{LWL}{Lichtwellenleiter}
\newacronym{DC/DC-Wandler}{DC/DC-Wandler}{Gleichspannungswandler}
\newacronym{JFET}{JFET}{junction field effect transistor}
\newacronym{DMOS}{DMOS}{double-diffused metal-oxide semiconductor}
\newacronym{PT}{PT}{Punch-Through}
\newacronym{SOA}{SOA}{Safe Operating Area}
\newacronym{PCB}{PCB}{Printed circuit board}
%\newacronym{SMD}{SMD}{Surface-mounted device}
\newglossaryentry{SMD}{
	name={SMD},
	description={Surface-mounted device},
	first={\glsentrydesc{SMD} (\glsentrytext{SMD})},
	plural={SMD},
	symbol={SMD},
	descriptionplural={Surface-mounted devices},
	firstplural={\glsentrydescplural{SMD} (\glsentryplural{SMD})}
}
\newacronym{THT}{THT}{Through-hole technology}
\newacronym{IC}{IC}{Integrated circuit}
\newacronym{Si}{Si}{Silizium}
\newacronym{C}{C}{Kohlenstoff}
\newacronym{SiC}{SiC}{Siliziumkarbid}
\newacronym{GaN}{GaN}{Galliumnitrid}
\newacronym{AlGaN}{AlGaN}{Aluminiumgalliumnitrid}
\newacronym{AlN}{AlN}{Aluminiumnitrid}
\newacronym{2DEG}{2DEG}{Zweidimensionales Elektronengas}
\newacronym{DC}{DC}{Gleichstrom (Direct Current)}
\newglossaryentry{AC}{
	name={AC},
	description={Wechselstrom (Alternating Current)},
	first={\glsentrydesc{AC} (\glsentrytext{AC})},
	plural={ACs},
	descriptionplural={Wechselströme (Alternating Currents)},
	firstplural={\glsentrydescplural{AC} (\glsentryplural{AC})}
}
\newacronym{TA}{TA}{Taktperiode des Umrichtersystems}
\newacronym{ETI}{ETI}{Elektrotechnisches Institut}

\newacronym{KIT}{KIT}{Karlsruher Institut für Technologie}
\newglossaryentry{EPSR}{
	name={EPSR},
	description={Einplatinenstromrichter},
	first={\glsentrydesc{EPSR} (\glsentrytext{EPSR})},
	plural={EPSR},
	descriptionplural={Einplatinenstromrichter},
	firstplural={\glsentrydescplural{EPSR} (\glsentryplural{EPSR})}
}
\newacronym{EMV}{EMV}{Elektromagnetischen Verträglichkeit}
\newacronym{FPGA}{FPGA}{Field Programmable Gate Array}
\newacronym{SoC}{SoC}{System on Chip}
\newacronym{DBS}{DBS}{Drehstrombrückenschaltung}
\newglossaryentry{DSP}{
	name={DSP},
	description={digitaler Signalprozessor},
	first={\glsentrydesc{DSP} (\glsentrytext{DSP})},
	plural={DSP's},
	descriptionplural={digitale Signalprozessoren},
	firstplural={\glsentrydescplural{DSP} (\glsentryplural{DSP})}
} 
\newacronym{ETI-Bus}{ETI-Bus}{kommunikationsbus des DSP-Systems}
\newacronym{USB}{USB}{Universal Serial Bus}
\newacronym{MMC}{MMC}{Modularer-Multilevel-Umrichter}
\newacronym{M3C}{M3C}{Modularer-Multilevel-Matrix-Umrichter}
\newacronym{ESB}{ESB}{Ersatzschaltbild}
\newacronym{3AC}{3AC}{dreiphasiger Wechselstrom}
\newacronym{AFE}{AFE}{Active Front End}
\newglossaryentry{TSS}{
	name={TSS},
	description={Tiefsetzsteller},
	first={\glsentrydesc{TSS} (\glsentrytext{TSS})},
	plural={TSS},
	descriptionplural={Tiefsetzsteller},
	firstplural={\glsentrydescplural{TSS} (\glsentryplural{TSS})}
} 

\newglossaryentry{DUT}{
	name={DUT},
	description={Device Under Test},
	first={\glsentrydesc{DUT} (\glsentrytext{DUT})}
} 
\newacronym{MAC}{MAC}{Multiply-Accumulate}

\newacronym{AD}{AD}{Analog Digital}
\newacronym{ADC}{ADC}{Analog Digital Wandler}
\newacronym{HMK}{HMK}{Hochleistungsmodulatorkarte}
\newacronym{GM}{GM}{Gleichstrommaschine}
\newacronym{ASM}{ASM}{Asynchronmaschine}
\newacronym{PLL}{PLL}{Phase Look Loop}
\newacronym{DG}{DG}{Diodengleichrichter}
\newacronym{PWM}{PWM}{Pulsweitenmodulation (Puls Wide Modulation)}
\newacronym{HIL}{HIL}{Hardware-in-the-Loop}
\newacronym{PHIL}{PHIL}{Power-Hardware-in-the-Loop}
\newacronym{SIL}{SIL}{Software-in-the-Loop}
\newglossaryentry{HB}{
	name={HB},
	description={Halbbrücke},
	first={\glsentrydesc{HB} (\glsentrytext{HB})},
	plural={HBs},
	descriptionplural={Halbbrücken},
	firstplural={\glsentrydescplural{HB} (\glsentryplural{HB})}
} 
\newglossaryentry{LED}{
	name={LED},
	description={light emitting diode},
	first={\glsentrydesc{LED} (\glsentrytext{LED})},
	plural={LED's},
	descriptionplural={light emitting diodes},
	firstplural={\glsentrydescplural{LED} (\glsentryplural{LED})}
} 
\newglossaryentry{EEPROM}{
	name={EEPROM},
	description={electrically erasable programmable read-only memory},
	first={\glsentrydesc{EEPROM} (\glsentrytext{EEPROM})},
	plural={EEPROM's},
	descriptionplural={electrically erasable programmable read-only memories},
	firstplural={\glsentrydescplural{EEPROM} (\glsentryplural{EEPROM})}
} 
\newacronym{FB}{FB}{Full Bridge - Vollbrücken}
\newacronym{PCC}{PCC}{Point of common coupling - Netzanschlusspunkt}
\newacronym{BLDC}{BLDC}{Bürstenloser Gleichstrommotor}
\newacronym{PSM}{PSM}{Permanenterregte Synchronmaschiene}
\newglossaryentry{MOSFET}{
	name={MOSFET},
	description={Metall Oxid Halbleiter Feldeffekttransistor},
	first={\glsentrydesc{MOSFET} (\glsentrytext{MOSFET})},
	plural={MOSFETs},
	descriptionplural={Metall Oxid Halbleiter Feldeffekttransistoren},
	firstplural={\glsentrydescplural{MOSFET} (\glsentryplural{MOSFET})}
}
\newacronym{OPV}{OPV}{Operationsverstärker}
\newacronym{VHDL}{VHDL}{Very High Speed Integrated Circuit Hardware Description Language}
\newacronym{IGBT}{IGBT}{Insulated Gate Bipolar Transistor}
\newacronym{usw.}{usw.}{und so weiter}
\newacronym{etc.}{etc.}{et cetera}
\newacronym{z.B.}{z.B.}{zum Beispiel}
\newglossaryentry{FAM}{
	name={FAM},
	description={Feldorientierte Regelung der Asynchronmaschine},
	first={\glsentrydesc{FAM} (\glsentrytext{FAM})},
	user1={Feldorientierten Regelung der Asynchronmaschine (\glsentrytext{FAM})}
} 

%--------------------------------------------------------------------
% die Auflistung für das Symbolverzeichnis
%--------------------------------------------------------------------


%% Allgemeine Größen Leistungselektronik

\newglossaryentry{sym:tau}{
	name={\ensuremath{\tau_\mathrm{G}}},
	user6={\glsentrydesc{sym:tau} \glsentrytext{sym:tau}},
	description={Zeitkonstante Tau},
	sort={tau}, 
	type={symbolslist}
}
\newglossaryentry{sym:ta}{
	name={\ensuremath{t_\mathrm{a}}},
	user6={\glsentrydesc{sym:ta} \glsentrytext{sym:ta}},
	description={Einschaltzeit des ersten Pulses eines Doppelpulsversuchs},
	sort={ta}, 
	type={symbolslist}
}
\newglossaryentry{sym:n-}{
	name={\ensuremath{n^\mathrm{-}}},
	user6={\glsentrydesc{sym:n-} \glsentrytext{sym:n-}},
	description={niedrig dotiertes Donator Halbleitermaterial},
	sort={n-}, 
	type={symbolslist}
}
\newglossaryentry{sym:SiO2}{
	name={\ensuremath{SiO_\mathrm{2}}},
	user6={\glsentrydesc{sym:SiO2} \glsentrytext{sym:SiO2}},
	description={Siliziumdioxid},
	sort={SiO2}, 
	type={symbolslist}
}

\newglossaryentry{sym:n+}{
	name={\ensuremath{n^\mathrm{+}}},
	user6={\glsentrydesc{sym:n+} \glsentrytext{sym:n+}},
	description={hoch dotiertes Donator Halbleitermaterial},
	sort={n+}, 
	type={symbolslist}
}

\newglossaryentry{sym:p+}{
	name={\ensuremath{p^\mathrm{+}}},
	user6={\glsentrydesc{sym:p+} \glsentrytext{sym:p+}},
	description={hoch dotiertes Akzeptor Halbleitermaterial},
	sort={p+}, 
	type={symbolslist}
}

\newglossaryentry{sym:p}{
	name={\ensuremath{p}},
	user6={\glsentrydesc{sym:p} \glsentrytext{sym:p}},
	description={p-dotiertes Akzeptor Halbleitermaterial},
	sort={p}, 
	type={symbolslist}
}
\newglossaryentry{sym:n}{
	name={\ensuremath{n}},
	user6={\glsentrydesc{sym:n} \glsentrytext{sym:n}},
	description={n-dotiertes Donator Halbleitermaterial},
	sort={n}, 
	type={symbolslist}
}

\newglossaryentry{sym:p-}{
	name={\ensuremath{p^\mathrm{-}}},
	user6={\glsentrydesc{sym:p-} \glsentrytext{sym:p-}},
	description={niedrig dotiertes Akzeptor Halbleitermaterial},
	sort={p-}, 
	type={symbolslist}
}

\newglossaryentry{sym:cdc}{
	name={\ensuremath{C_\mathrm{DC}}},
	user6={\glsentrydesc{sym:cdc} \glsentrytext{sym:cdc}},
	description={Zwischenkreis-Kapazität},
	sort={cdc}, 
	type={symbolslist}
}


\newglossaryentry{sym:rg}{
	name={\ensuremath{R_\mathrm{G}}},
	user6={\glsentrydesc{sym:rg} \glsentrytext{sym:rg}},
	description={Gate-Widerstand},
	sort={rg}, 
	type={symbolslist}
}
\newglossaryentry{sym:llast}{
	name={\ensuremath{L_\mathrm{Last}}},
	user6={\glsentrydesc{sym:llast} \glsentrytext{sym:llast}},
	description={Last-Induktivität},
	sort={llast}, 
	type={symbolslist}
}

\newglossaryentry{sym:rgon}{
	name={\ensuremath{R_\mathrm{G,on}}},
	user6={\glsentrydesc{sym:rgon} \glsentrytext{sym:rgon}},
	description={Gate-Widerstand im Einschaltvorgang},
	sort={rgon}, 
	type={symbolslist}
}
\newglossaryentry{sym:rgoff}{
	name={\ensuremath{R_\mathrm{G,off}}},
	user6={\glsentrydesc{sym:rgoff} \glsentrytext{sym:rgoff}},
	description={Gate-Widerstand im Ausschaltvorgang},
	sort={rgoff}, 
	type={symbolslist}
}
\newglossaryentry{sym:rgint}{
	name={\ensuremath{R_\mathrm{G,int}}},
	user6={\glsentrydesc{sym:rgint} \glsentrytext{sym:rgint}},
	description={Interner Gate-Widerstand},
	sort={rgint}, 
	type={symbolslist}
}
\newglossaryentry{sym:rd}{
	name={\ensuremath{R_\mathrm{D}}},
	user6={\glsentrydesc{sym:rd} \glsentrytext{sym:rd}},
	description={Driftwiderstand des n- Gebiets},
	sort={rd}, 
	type={symbolslist}
}
\newglossaryentry{sym:eon}{
	name={\ensuremath{E_\mathrm{ON}}},
	user6={\glsentrydesc{sym:eon} \glsentrytext{sym:eon}},
	description={Einschaltenergie},
	sort={eon}, 
	type={symbolslist}
}
\newglossaryentry{sym:eoff}{
	name={\ensuremath{E_\mathrm{OFF}}},
	user6={\glsentrydesc{sym:eoff} \glsentrytext{sym:eoff}},
	description={Ausschaltenergie},
	sort={eoff}, 
	type={symbolslist}
}
\newglossaryentry{sym:lsigma}{
	name={\ensuremath{L_\mathrm{\sigma}}},
	user6={\glsentrydesc{sym:lsigma} \glsentrytext{sym:lsigma}},
	description={parasitäre Induktivität im Kommutierungskreis},
	sort={lsigma}, 
	type={symbolslist}
}
\newglossaryentry{sym:udc}{
	name={\ensuremath{U_\mathrm{DC}}},
	user6={\glsentrydesc{sym:udc} \glsentrytext{sym:udc}},
	description={Zwischenkreisspannung},
	sort={udc}, 
	type={symbolslist}
}
\newglossaryentry{sym:mue_p}{
	name={\ensuremath{\mu_\mathrm{p}}},
	description={Beweglichkeit der Löcher},
	sort={mue_p}, 
	type={symbolslist}
}
\newglossaryentry{sym:mue_n}{
	name={\ensuremath{\mu_\mathrm{n}}},
	description={Beweglichkeit der Elektronen},
	sort={mue_n}, 
	type={symbolslist}
}
\newglossaryentry{sym:irr}{
	name={\ensuremath{I_\mathrm{rr}}},
	user6={\glsentrydesc{sym:irr} \glsentrytext{sym:irr}},
	description={Diodenrückstrom},
	sort={irr}, 
	type={symbolslist}
}
\newglossaryentry{sym:id}{
	name={\ensuremath{I_\mathrm{D}}},
	user6={\glsentrydesc{sym:id} \glsentrytext{sym:id}},
	description={Drainstrom},
	sort={id}, 
	type={symbolslist}
}
\newglossaryentry{sym:qrr}{
	name={\ensuremath{Q_\mathrm{rr}}},
	user6={\glsentrydesc{sym:qrr} \glsentrytext{sym:qrr}},
	description={Reverse recovery Speicherladung einer Diode},
	sort={qrr}, 
	type={symbolslist}
}
\newglossaryentry{sym:ilast}{
	name={\ensuremath{I_\mathrm{Last}}},
	user6={\glsentrydesc{sym:ilast} \glsentrytext{sym:ilast}},
	description={Laststrom},
	sort={irrm}, 
	type={symbolslist}
}
\newglossaryentry{sym:ig}{
	name={\ensuremath{I_\mathrm{G}}},
	user6={\glsentrydesc{sym:ig} \glsentrytext{sym:ig}},
	description={Gatestrom},
	sort={ig}, 
	type={symbolslist}
}
\newglossaryentry{sym:irrm}{
	name={\ensuremath{I_\mathrm{rrm}}},
	user6={\glsentrydesc{sym:irrm} \glsentrytext{sym:irrm}},
	description={Rückstromspitze},
	sort={irrm}, 
	type={symbolslist}
}
\newglossaryentry{sym:irms}{
	name={\ensuremath{I_\mathrm{RMS}}},
	user6={\glsentrydesc{sym:irms} \glsentrytext{sym:irms}},
	description={Effektivwert des Stroms},
	sort={irrm}, 
	type={symbolslist}
}
\newglossaryentry{sym:ugon}{
	name={\ensuremath{U_\mathrm{G,on}}},
	user6={\glsentrydesc{sym:ugon} \glsentrytext{sym:ugon}},
	description={Ausgangsspannungpegel eines Gatetreiber zum Einschalten eines Transistors},
	sort={ugon}, 
	type={symbolslist}	
}
\newglossaryentry{sym:uuvlods}{
	name={\ensuremath{U_\mathrm{UVLO,DS}}},
	user6={\glsentrydesc{sym:uuvlods} \glsentrytext{sym:uuvlods}},
	description={angegebener UVLO Schwellwert im Datenblatt eines Gatetreibers},
	sort={uuvlods}, 
	type={symbolslist}	
}

\newglossaryentry{sym:ugoff}{
	name={\ensuremath{U_\mathrm{G,off}}},
	user6={\glsentrydesc{sym:ugoff} \glsentrytext{sym:ugoff}},
	description={Ausgangsspannungpegel eines Gatetreiber zum Ausschalten eines Transistors},
	sort={ugoff}, 
	type={symbolslist}	
}
\newglossaryentry{sym:ugtreiber}{
	name={\ensuremath{u_\mathrm{G,Treiber}}},
	user6={\glsentrydesc{sym:ugtreiber} \glsentrytext{sym:ugtreiber}},
	description={Steuerspannung am Ausgang eines Gatetreibers},
	sort={ugtreiber}, 
	type={symbolslist}	
}
\newglossaryentry{sym:deltauls}{
	name={\ensuremath{\Delta U_\mathrm{LS}}},
	user6={\glsentrydesc{sym:deltauls} \glsentrytext{sym:deltauls}},
	description={Induzierte Gegespannung im Ausschaltvorgang},
	sort={deltauls}, 
	type={symbolslist}	
}
\newglossaryentry{sym:deltaul}{
	name={\ensuremath{\Delta U_\mathrm{L}}},
	user6={\glsentrydesc{sym:deltaus} \glsentrytext{sym:deltaus}},
	description={Induzierter Spannungsabfall im Einschaltvorgang},
	sort={deltaus}, 
	type={symbolslist}	
}
\newglossaryentry{sym:td}{
	name={\ensuremath{t_\mathrm{D}}},
	user6={\glsentrydesc{sym:td} \glsentrytext{sym:td}},
	description={Verriegelungszeit},
	sort={td}, 
	type={symbolslist}	
}
%% IGBT Kapazitäten und Spannungen

\newglossaryentry{sym:ices}{
	name={\ensuremath{I_\mathrm{CES}}},
	user6={\glsentrydesc{sym:ices} \glsentrytext{sym:ices}},
	description={Kollektor Reststrom},
	sort={ices}, 
	type={symbolslist}
}
\newglossaryentry{sym:rw}{
	name={\ensuremath{R_\mathrm{W}}},
	user6={\glsentrydesc{sym:rw} \glsentrytext{sym:rw}},
	description={Wannenwiderstand},
	sort={rw}, 
	type={symbolslist}
}
\newglossaryentry{sym:ic}{
	name={\ensuremath{I_\mathrm{C}}},
	user6={\glsentrydesc{sym:ic} \glsentrytext{sym:ic}},
	description={Kollektor-Strom},
	sort={ic}, 
	type={symbolslist}
}
\newglossaryentry{sym:cge}{
	name={\ensuremath{C_\mathrm{GE}}},
	user6={\glsentrydesc{sym:cge} \glsentrytext{sym:cge}},
	description={Gate-Emitter-Kapazität},
	sort={cge}, 
	type={symbolslist}
}
\newglossaryentry{sym:cce}{
	name={\ensuremath{C_\mathrm{CE}}},
	user6={\glsentrydesc{sym:cce} \glsentrytext{sym:cce}},
	description={Kollektor-Emitter-Kapazität},
	sort={cce}, 
	type={symbolslist}
}
\newglossaryentry{sym:cgc}{
	name={\ensuremath{C_\mathrm{GC}}},
	user6={\glsentrydesc{sym:cgc} \glsentrytext{sym:cgc}},
	description={Gate-Kollektor-Kapazität},
	sort={cge}, 
	type={symbolslist}
}

\newglossaryentry{sym:ciss}{
	name={\ensuremath{C_\mathrm{iss}}},
	user6={\glsentrydesc{sym:ciss} \glsentrytext{sym:ciss}},
	description={Eingangskapazität},
	sort={ciss}, 
	type={symbolslist}
}	
\newglossaryentry{sym:coss}{
	name={\ensuremath{C_\mathrm{oss}}},
	user6={\glsentrydesc{sym:coss} \glsentrytext{sym:coss}},
	description={Ausgangskapazität},
	sort={coss}, 
	type={symbolslist}
}
\newglossaryentry{sym:crss}{
	name={\ensuremath{C_\mathrm{rss}}},
	user6={\glsentrydesc{sym:crss} \glsentrytext{sym:crss}},
	description={Rückwirkungskapazität},
	sort={crss}, 
	type={symbolslist}
}
\newglossaryentry{sym:uge}{
	name={\ensuremath{U_\mathrm{GE}}},
	user6={\glsentrydesc{sym:uge} \glsentrytext{sym:uge}},
	description={Gate-Emitter-Spannung},
	sort={uge}, 
	type={symbolslist}
}
\newglossaryentry{sym:ugeth}{
	name={\ensuremath{U_\mathrm{GE,th}}},
	user6={\glsentrydesc{sym:ugeth} \glsentrytext{sym:ugeth}},
	description={Threshold-Spannung eines IGBT},
	sort={ugeth}, 
	type={symbolslist}	
}
\newglossaryentry{sym:ugeon}{
	name={\ensuremath{U_\mathrm{GE,on}}},
	user6={\glsentrydesc{sym:ugeon} \glsentrytext{sym:ugeon}},
	description={Gate-Emitter Spannung im eingeschalteten Zustand},
	sort={ugeon}, 
	type={symbolslist}	
}
\newglossaryentry{sym:ugeoff}{
	name={\ensuremath{U_\mathrm{GE,off}}},
	user6={\glsentrydesc{sym:ugeoff} \glsentrytext{sym:ugeoff}},
	description={Gate-Emitter Spannung im ausgeschalteten Zustand},
	sort={ugeon}, 
	type={symbolslist}	
}
\newglossaryentry{sym:uce}{
	name={\ensuremath{U_\mathrm{CE}}},
	user6={\glsentrydesc{sym:uce} \glsentrytext{sym:uce}},
	description={Kollektor-Emitter-Spannung},
	sort={uce}, 
	type={symbolslist}
}
\newglossaryentry{sym:ucesat}{
	name={\ensuremath{U_\mathrm{CE(sat)}}},
	user6={\glsentrydesc{sym:ucesat} \glsentrytext{sym:ucesat}},
	description={Kollektor-Emitter-Spannung in Sättigung},
	sort={ucesat}, 
	type={symbolslist}
}
\newglossaryentry{sym:uth}{
	name={\ensuremath{U_\mathrm{th}}},
	user6={\glsentrydesc{sym:uth} \glsentrytext{sym:uth}},
	description={Threshold-Spannung},
	sort={uth}, 
	type={symbolslist}
}


%% Mosfet Kapazitäten und Spannungen
\newglossaryentry{sym:uuvlogan}{
	name={\ensuremath{U_\mathrm{UVLO,GaN}}},
	user6={\glsentrydesc{sym:uuvlogan} \glsentrytext{sym:uuvlogan}},
	description={Effektive undervoltage-lockout Schwelle des GaN-Treibers},
	sort={uuvlogan}, 
	type={symbolslist}
}
\newglossaryentry{sym:idss}{
	name={\ensuremath{I_\mathrm{DSS}}},
	user6={\glsentrydesc{sym:idss} \glsentrytext{sym:idss}},
	description={Drain Reststrom},
	sort={idss}, 
	type={symbolslist}
}
\newglossaryentry{sym:ugsoff}{
	name={\ensuremath{U_\mathrm{GS,off}}},
	user6={\glsentrydesc{sym:ugsoff} \glsentrytext{sym:ugsoff}},
	description={Gate-Source Spannung im ausgeschalteten Zustand },
	sort={ugsoff}, 
	type={symbolslist}
}

\newglossaryentry{sym:ugson}{
	name={\ensuremath{U_\mathrm{GS,on}}},
	user6={\glsentrydesc{sym:ugson} \glsentrytext{sym:ugson}},
	description={Gate-Source Spannung im eingeschalteten Zustand },
	sort={ugson}, 
	type={symbolslist}
}
\newglossaryentry{sym:ufnull}{
	name={\ensuremath{U_\mathrm{F,0}}},
	user6={\glsentrydesc{sym:ufnull} \glsentrytext{sym:ufnull}},
	description={Flussspannung der parasitären MOSFET body-Diode},
	sort={ufnull}, 
	type={symbolslist}
}
\newglossaryentry{sym:udson}{
	name={\ensuremath{U_\mathrm{DS,on}}},
	user6={\glsentrydesc{sym:udson} \glsentrytext{sym:udson}},
	description={Drain-Source Spannung im eingeschalteten Zustand eines MOSFET},
	sort={udson}, 
	type={symbolslist}
}

\newglossaryentry{sym:rdson}{
	name={\ensuremath{R_\mathrm{DS(on)}}},
	user6={\glsentrydesc{sym:rdson} \glsentrytext{sym:rdson}},
	description={Durchlasswiderstand},
	sort={rdson}, 
	type={symbolslist}
}

\newglossaryentry{sym:cgd}{
	name={\ensuremath{C_\mathrm{GD}}},
	user6={\glsentrydesc{sym:cgd} \glsentrytext{sym:cgd}},
	description={Gate-Drain-Kapazität},
	sort={cgd}, 
	type={symbolslist}
}
\newglossaryentry{sym:cgs}{
	name={\ensuremath{C_\mathrm{GS}}},
	user6={\glsentrydesc{sym:cgs} \glsentrytext{sym:cgs}},
	description={Gate-Source-Kapazität},
	sort={cgs}, 
	type={symbolslist}
}

\newglossaryentry{sym:cds}{
	name={\ensuremath{C_\mathrm{DS}}},
	user6={\glsentrydesc{sym:cds} \glsentrytext{sym:cds}},
	description={Drain-Source-Kapazität},
	sort={cds}, 
	type={symbolslist}
}
\newglossaryentry{sym:didt}{
	name={\ensuremath{\mathrm{d}i/\mathrm{d}t}},
	user6={\glsentrydesc{sym:didt} \glsentrytext{sym:didt}},
	description={Stromgradient},
	sort={didt}, 
	type={symbolslist}
}
\newglossaryentry{sym:dudt}{
	name={\ensuremath{\mathrm{d}u/\mathrm{d}t}},
	user6={\glsentrydesc{sym:dudt} \glsentrytext{sym:dudt}},
	description={Spannungsgradient},
	sort={dudt}, 
	type={symbolslist}
}
\newglossaryentry{sym:ugs}{
	name={\ensuremath{U_\mathrm{GS}}},
	user6={\glsentrydesc{sym:ugs} \glsentrytext{sym:ugs}},
	description={Gate-Source-Spannung},
	sort={ugs}, 
	type={symbolslist}	
}
\newglossaryentry{sym:ugsth}{
	name={\ensuremath{U_\mathrm{GS,th}}},
	user6={\glsentrydesc{sym:ugsth} \glsentrytext{sym:ugsth}},
	description={Threshold-Spannung eines MOSFET},
	sort={ugsth}, 
	type={symbolslist}
}
\newglossaryentry{sym:uds}{
	name={\ensuremath{U_\mathrm{DS}}},
	user6={\glsentrydesc{sym:uds} \glsentrytext{sym:uds}},
	description={Drain-Source-Spannung},
	sort={uds}, 
	type={symbolslist}
}

%% Allgemeine Größen
\newglossaryentry{sym:pollparzahl}{
	name={\ensuremath{p}},
	user6={\glsentrydesc{sym:pollparzahl} \glsentrytext{sym:pollparzahl}},
	description={Polpaarzahl der Asynchronmaschine},
	sort={polPar}, 
	type={symbolslist}
}

\newglossaryentry{sym:inneresMoment}{
	name={\ensuremath{M_\mathrm{i}}},
	user6={\glsentrydesc{sym:inneresMoment} \glsentrytext{sym:inneresMoment}},
	description={inneres Moment der Asynchronmaschine},
	sort={Minner}, 
	type={symbolslist}
}

\newglossaryentry{sym:lastMoment}{
	name={\ensuremath{M_\mathrm{L}}},
	user6={\glsentrydesc{sym:lastMoment} \glsentrytext{sym:lastMoment}},
	description={Lastmoment der Asynchronmaschine},
	sort={Mlast}, 
	type={symbolslist}
}

\newglossaryentry{sym:reibMoment}{
	name={\ensuremath{M_\mathrm{Reib}}},
	user6={\glsentrydesc{sym:reibMoment} \glsentrytext{sym:reibMoment}},
	description={Reibmoment der Asynchronmaschine},
	sort={Mreib}, 
	type={symbolslist}
}

\newglossaryentry{sym:traegheitsmoment}{
	name={\ensuremath{J_\mathrm{Sys}}},
	user6={\glsentrydesc{sym:traegheitsmoment} \glsentrytext{sym:traegheitsmoment}},
	description={Trägheitsmoment des mechanischen Systems},
	sort={Jsys}, 
	type={symbolslist}
}


\newglossaryentry{sym:indukhaupt}{
	name={\ensuremath{L_\mathrm{h}}},
	user6={\glsentrydesc{sym:indukhaupt} \glsentrytext{sym:indukhaupt}},
	description={Hauptinduktivität},
	sort={Lh}, 
	type={symbolslist}
}
\newglossaryentry{sym:magnetisierungsstrom}{
	name={\ensuremath{i_\mathrm{\mu}}},
	user6={\glsentrydesc{sym:magnetisierungsstrom} \glsentrytext{sym:magnetisierungsstrom}},
	description={Magnetisierungsstrom der Asynchronmaschine},
	sort={imagnetstrom}, 
	type={symbolslist}
}

\newglossaryentry{sym:omegamech}{
	name={\ensuremath{\Omega_\mathrm{Mech}}},
	user2= \ensuremath{\Omega_\mathrm{Mech}(t)}, 					% Zeitverlauf der mechanischen Größe 
	user3= \ensuremath{\omega_\mathrm{Mech,0}},
	user4= \ensuremath{\omega_\mathrm{Mech,1}},
	user6={\glsentrydesc{sym:ommegamech} \glsentrytext{sym:ommegamech}},
	description={mechanische Drehfrequenz des Rotors der Asynchronmaschine},
	sort={Omegastator}, 
	type={symbolslist}
}

%% Alle größen für die Statorseite

\newglossaryentry{sym:statorwiderstand}{
	name={\ensuremath{R_\mathrm{S}}},
	user6={\glsentrydesc{sym:statorwiderstand} \glsentrytext{sym:statorwiderstand}},
	description={Statorwiderstand 123},
	sort={Rs}, 
	type={symbolslist}
}
\newglossaryentry{sym:statorspannungkompl}{
	name={\ensuremath{\uline{u}_\mathrm{S}}},
	user6={\glsentrydesc{sym:statorspannungkompl} \glsentrytext{sym:statorspannungkompl}},
	description={komplexe Statorspannung},
	sort={us}, 
	type={symbolslist}
}
\newglossaryentry{sym:statorstromkompl}{
	name={\ensuremath{\uline{i}_\mathrm{S}}},
	user6={\glsentrydesc{sym:statorstromkompl} \glsentrytext{sym:statorstromkompl}},
	description={komplexer Statorstrom},
	sort={is}, 
	type={symbolslist}
}
\newglossaryentry{sym:statorflusskompl}{
	name={\ensuremath{\uline{\Psi}_\mathrm{S}}},			
	user1={\ensuremath{\uline{\dot{\Psi}}_\mathrm{S}}},			% 1. Ableitung
	user6={\glsentrydesc{sym:statorflusskompl} \glsentrytext{sym:statorflusskompl}},
	description={komplexer Statorfluss},
	sort={Psis}, 
	type={symbolslist}
}
\newglossaryentry{sym:indukstatorstreu}{
	name={\ensuremath{L_\mathrm{S\sigma}}},
	user6={\glsentrydesc{sym:indukstatorstreu} \glsentrytext{sym:indukstatorstreu}},
	description={Statorstreuinduktivität},
	sort={Lssigma}, 
	type={symbolslist}
}

\newglossaryentry{sym:gammarstator}{
	name={\ensuremath{\gamma_\mathrm{S}}},
	user2={\ensuremath{\dot{\gamma}_\mathrm{S}}},	 % Ableitung
	user6={\glsentrydesc{sym:gammarstator} \glsentrytext{sym:gammarstator}},
	description={Winkel des Statorfluss der Asynchronmaschine},
	sort={gammarstator}, 
	type={symbolslist}
}

\newglossaryentry{sym:ommegastator}{
	name={\ensuremath{\omega_\mathrm{S}}},
	user6={\glsentrydesc{sym:ommegastator} \glsentrytext{sym:ommegastator}},
	description={Speißefrequenz des Statorflusses der Asynchronmaschine},
	sort={Omegastator}, 
	type={symbolslist}
}

\newglossaryentry{sym:ommegasyn}{
	name={\ensuremath{\omega_\mathrm{Syn}}},
	user6={\glsentrydesc{sym:ommegastator} \glsentrytext{sym:ommegastator}},
	description={synchrone Kreisfrequenz},
	sort={Omegasyn}, 
	type={symbolslist}
}

\newglossaryentry{sym:ommegaeck}{
	name={\ensuremath{\omega_\mathrm{Eck}}},
	user6={\glsentrydesc{sym:ommegastator} \glsentrytext{sym:ommegastator}},
	description={Eckkreisfrequenz (beginn des Feldschwächebereich)},
	sort={Omegaeck}, 
	type={symbolslist}
}

%% Alle größen für die Rotorseite

\newglossaryentry{sym:rotorwiderstand}{
	name={\ensuremath{R_\mathrm{R}}},
	user2={\ensuremath{R_\mathrm{R}^\prime}},						% Transformiert
	user6={\glsentrydesc{sym:rotorwiderstand} \glsentrytext{sym:rotorwiderstand}},
	description={Rotorwiderstand},
	sort={Rr}, 
	type={symbolslist}
}
\newglossaryentry{sym:rotorspannungkompl}{
	name={\ensuremath{\uline{u}_\mathrm{R}}},
	user6={\glsentrydesc{sym:rotorspannungkompl} \glsentrytext{sym:rotorspannungkompl}},
	description={komplexe Rotorspannung},
	sort={ur}, 
	type={symbolslist}
}
\newglossaryentry{sym:rotorstromkompl}{
	name={\ensuremath{\uline{i}_\mathrm{R}}},
	user2={\ensuremath{\uline{i}_\mathrm{R}^\prime}},				% Transformiert
	user3={\ensuremath{\uline{i}_\mathrm{R}^{\prime*}}},				% konjugiert komplex
	user6={\glsentrydesc{sym:rotorstromkompl} \glsentrytext{sym:rotorstromkompl}},
	description={komplexer Rotorstrom},
	sort={ir}, 
	type={symbolslist}
}
\newglossaryentry{sym:rotorflusskompl}{
	name={\ensuremath{\uline{\Psi}_\mathrm{R}}},			
	user1={\ensuremath{\uline{\dot{\Psi}}_\mathrm{R}}},				% 1. Ableitung
	user2={\ensuremath{\uline{\Psi}_\mathrm{R}^\prime}},			% Transformiert
	user3={\ensuremath{\uline{\dot{\Psi}}_\mathrm{R}^\prime}},		% Transformiert, 1. Ableitung
	user6={\glsentrydesc{sym:rotorflusskompl} \glsentrytext{sym:rotorflusskompl}},
	description={komplexer Rotorfluss},
	sort={Psir}, 
	type={symbolslist}
}
\newglossaryentry{sym:gammarrotor}{
	name={\ensuremath{\gamma_\mathrm{R}}},
	user2={\ensuremath{\dot{\gamma}_\mathrm{R}}},	 % Ableitung
	user6={\glsentrydesc{sym:gammarrotor} \glsentrytext{sym:gammarrotor}},
	description={Winkel des Rotorfluss der Asynchronmaschine},
	sort={gammarrotor}, 
	type={symbolslist}
}

\newglossaryentry{sym:indukrotorstreu}{
	name={\ensuremath{L_\mathrm{R\sigma}}},
	user2={\ensuremath{L_\mathrm{R\sigma}^\prime}},			% Transformiert
	user6={\glsentrydesc{sym:indukrotorstreu} \glsentrytext{sym:indukrotorstreu}},
	description={Rotorstreuinduktivität},
	sort={Lrsigma}, 
	type={symbolslist}
}

\newglossaryentry{sym:ommegarotor}{
	name={\ensuremath{\omega_\mathrm{R}}},
	user6={\glsentrydesc{sym:ommegarotor} \glsentrytext{sym:ommegarotor}},
	description={Speißefrequenz des Rotorflusses der Asynchronmaschine},
	sort={Omegastator}, 
	type={symbolslist}
}